%----------------------------------------------------------------------------------------
%	PACKAGES AND OTHER DOCUMENT CONFIGURATIONS
%----------------------------------------------------------------------------------------

\documentclass[
	11pt, % Set the default font size, options include: 8pt, 9pt, 10pt, 11pt, 12pt, 14pt, 17pt, 20pt
	%t, % Uncomment to vertically align all slide content to the top of the slide, rather than the default centered
	aspectratio=169, % Uncomment to set the aspect ratio to a 16:9 ratio which matches the aspect ratio of 1080p and 4K screens and projectors
	envcountsect,
	xcolor=dvipsnames,
	handout
]{beamer}

\usepackage{PresentationTemplate}

\graphicspath{{Images/}{./}} % Specifies where to look for included images (trailing slash required)
%----------------------------------------------------------------------------------------
%	PRESENTATION INFORMATION
%----------------------------------------------------------------------------------------

\title[CVE101 - Civil Engineering Orientation]{Lecture 1} % The short title in the optional parameter appears at the bottom of every slide, the full title in the main parameter is only on the title page

\subtitle{What is Civil Engineering} % Presentation subtitle, remove this command if a subtitle isn't required

\author[Nikol O. Telen]{Nikol O. Telen} % Presenter name(s), the optional parameter can contain a shortened version to appear on the bottom of every slide, while the main parameter will appear on the title slide

\institute[MSU-GenSan, Civil Engineering Department
]{{\scriptsize Mindanao State University - General Santos}\\\textbf{CIVIL ENGINEERING DEPARTMENT}} % Your institution, the optional parameter can be used for the institution shorthand and will appear on the bottom of every slide after author names, while the required parameter is used on the title slide and can include your email address or additional information on separate lines

\date{\today} % Presentation date or conference/meeting name, the optional parameter can contain a shortened version to appear on the bottom of every slide, while the required parameter value is output to the title slide

\titlegraphic{
	\includegraphics[width=1cm]{MSU.jpg}
	\hspace{0.2cm}
	\includegraphics[width=1cm]{CEdept.png}
} 

\institute[MSU-GenSan, Civil Engineering Department]{{\scriptsize Mindanao State University - General Santos}\\\textbf{CIVIL ENGINEERING DEPARTMENT}}
\date{\today} 
\subject{CVE101 - Civil Engineering Orientation}

%----------------------------------------------------------------------------------------

\begin{document}

%----------------------------------------------------------------------------------------
%	TITLE SLIDE
%----------------------------------------------------------------------------------------

\begin{frame}
	\titlepage % Output the title slide, automatically created using the text entered in the PRESENTATION INFORMATION block above
\end{frame}

%----------------------------------------------------------------------------------------
%	TABLE OF CONTENTS SLIDE
%----------------------------------------------------------------------------------------

% The table of contents outputs the sections and subsections that appear in your presentation, specified with the standard \section and \subsection commands. You may either display all sections and subsections on one slide with \tableofcontents, or display each section at a time on subsequent slides with \tableofcontents[pausesections]. The latter is useful if you want to step through each section and mention what you will discuss.

\begin{frame}
	\frametitle{Lecture Overview} % Slide title, remove this command for no title
	
	\tableofcontents % Output the table of contents (all sections on one slide)
	%\tableofcontents[pausesections] % Output the table of contents (break sections up across separate slides)
\end{frame}

%----------------------------------------------------------------------------------------
%	PRESENTATION BODY SLIDES
%----------------------------------------------------------------------------------------

\section{Overview} % Sections are added in order to organize your presentation into discrete blocks, all sections and subsections are automatically output to the table of contents as an overview of the talk but NOT output in the presentation as separate slides	
\SectionPage
	\begin{frame}{Civil Engineering}
		Civil engineering is the oldest branch of engineering that deals with providing people with a livable built environment
		
		\bigskip

		It deals with people's everyday needs such as:
	
		\begin{itemize}
			\item<2-> food and water
			\item<3-> livable shelter
			\item<4-> ways to travel
			\item<5-> safety from disaster
		\end{itemize}

	\end{frame}

	\begin{frame}{Civil Engineering}
		Each basic need is related to at least one civil engineering technical area. 
	\end{frame}
\section{Civil Engineering Technical Areas}
\SectionPage
\subsection{Structural Engineering}
	\begin{frame}{Structural Engineering}
		Structural engineering is the technical specialty that deals with analysis and design of structures.
		
		\bigskip

		Structural design aims at providing a structure with sufficient level of resistance against these loads with minimum cost
		
	\end{frame}

	\begin{frame}{Structural Engineering}
		\framesubtitle{Technical Sub-Areas}
		It has several technical sub-areas, named according to the type of structure or to the type of load it carries.
		\begin{itemize}
			\item<2-> Earthquake Engineering
			\item<3-> Wind Engineering
			\item<4-> Structural Reliability
			\item<5-> Fire Engineering
			\item<6-> Bridge Engineering
			\item<7-> Dam Engineering
			\item<8-> Building Engineering
			\item<9-> Forensic Engineering
		\end{itemize}			
	\end{frame}

	\subsection{Geotechnical Engineering}
	\begin{frame}{Geotechnical Engineering}
		Geotechnical engineering is the technical specialty that deals with soil and rock as supporting materials for structures
		
		\bigskip

		It deals with the various foundation types that work between the structure and the ground.

		\bigskip

		It also deals with the stability of soil and rocks.
		
	\end{frame}

	\begin{frame}{Geotechnical Engineering}
		\framesubtitle{Technical Sub-Areas}
		Technical areas of study pertinent to geotechnical engineering includes:
		\begin{itemize}
			\item<2-> Engineering Geology
			\item<3-> Soil Mechanics
			\item<4-> Rock Mechanics
			\item<5-> Foundation Engineering
			\item<6-> Soil Improvement
			\item<7-> Tunnel Engineering
		\end{itemize}			
	\end{frame}

	\subsection{Water Resources Engineering}
	\begin{frame}{Water Resources Engineering}
		Water resource engineering is the specialty that deals with use of water in different human needs.
		
		\bigskip

		It includes finding and preservation of water sources, its means of delivery to a city or region, as well as studying its movement, and its erosive effects.
		
	\end{frame}

	\begin{frame}{Water Resources Engineering}
		\framesubtitle{Technical Sub-Areas}
		Its core and related areas of study include:
		\begin{itemize}
			\item<2-> Water Resources System Engineering
			\item<3-> Hydraulic Engineering
			\item<4-> Coastal Engineering
			\item<5-> Ocean Engineering
		\end{itemize}			
	\end{frame}

	\subsection{Transportation Engineering}
	\begin{frame}{Transportation Engineering}
		Transportation engineering deals with the efficient transport of people and goods.
		
		\bigskip

		Involved in the design of roads, harbors, and airport, as well as railways and traffic control.

	\end{frame}

	\begin{frame}{Transportation Engineering}
		\framesubtitle{Technical Sub-Areas}
		Technical specialties include:
		\begin{itemize}
			\item<2-> Transportation Planning
			\item<3-> Transportation System Engineering
			\item<4-> Highway Engineering
			\item<5-> Railway Engineering
			\item<6-> Port and Harbor Engineering
			\item<7-> Airport Engineering
			\item<8-> Traffic Engineering
		\end{itemize}			
	\end{frame}

	\subsection{Construction Engineering}
	\begin{frame}{Construction Engineering}
		Construction engineering deals with the execution of the plan designed on paper into physical reality.
		
		\bigskip

		It involves management of different resources required in the construction process.

	\end{frame}

	\begin{frame}{Construction Engineering}
		\framesubtitle{Technical Sub-Areas}
		Engineering and management aspects include:
		\begin{itemize}
			\item<2-> Construction Methods
			\item<3-> Electric and Mechanical Facilities
			\item<4-> Construction Machineries
			\item<5-> Financial and Cost management
			\item<6-> Contracts and Specifications
			\item<7-> Health and Safety Issues
			\item<8-> Legal Issues and Risk Management
		\end{itemize}			
	\end{frame}

\section{Related Disciplines}
\SectionPage
\subsection*{Geomatics (Survey Engineering)}
	\begin{frame}{Geomatics Engineering}
		A technical specialty that deals the measurement of and assembling spatial data of any natural or constructed objects in the earth.
		
		\bigskip

		It is involved in documentation of real estate, location of routes and points needed in construction, and collecting global data for resource analysis and utilization.
		
	\end{frame}

	\begin{frame}{Geomatics (Survey Engineering)}
		\framesubtitle{Technical Sub-Areas}
		Technical areas of study include:
		\begin{itemize}
			\item<2-> Plane Surveying
			\item<3-> Route Location
			\item<4-> Land Surveying
			\item<5-> Geodetic Surveying
			\item<6-> Aerial Photography and Satellite Imaging
			\item<7-> Geographical Information System
		\end{itemize}			
	\end{frame}

\subsection*{Environmental Engineering}
	\begin{frame}{Environmental Engineering}
		Environmental engineering involves the application engineering to protect human health and preserve the natural environment.
		
		\bigskip

		It relies on the knowledge of fundamental sciences of chemistry, biology, ecology, and health.
		
	\end{frame}

	\begin{frame}{Environmental Engineering}
		\framesubtitle{Technical Sub-Areas}
		Technical areas of study include:
		\begin{itemize}
			\item<2-> Water Treatment and Supply
			\item<3-> Wastewater Treatment and Disposal
			\item<4-> Solid Waste Disposal
		\end{itemize}			
	\end{frame}
\subsection*{Urban Planning}
	\begin{frame}{Urban Planning}
		Urban planning integrates land use planning, infrastructure planning, and
		 public policy for new developments  or  renewal  of  urbanized 
		 communities.  
		 
		 \bigskip 
		 
		 Successful  urban  planning  requires  the  application  of  the  knowledge developed in social, economic, architectural, and engineering studies
		
	\end{frame}

\section*{Other Related Disciplines}
\SectionPage
	\begin{frame}{Other Related Disciplines}
		The following disciplines interact or overlap with civil engineering:
		\begin{itemize}
			\item<2-> Applied Mechanics
			\item<3-> Architectural Engineering
			\item<4-> Agricultural Engineering
			\item<5-> Aerospace Engineering
			\item<6-> Biomedical Engineering
			\item<7-> Naval Architecture
		\end{itemize}
	\end{frame}

\section{Civil Engineering Processes}
\SectionPage
	\begin{frame}{Civil Engineering Processes}
		Any kind of civil engineering project involves four main phases:
		\begin{itemize}
			\item<2-> Planning
			\item<3-> Design
			\item<4-> Construction
			\item<5-> Maintenance/Operation
		\end{itemize}
	\end{frame}

\subsection*{Planning}
	\begin{frame}{Planning}
		\begin{itemize}
			\item<2-> The planning stage involves the study of its effects to private or public interest.
			\item<3-> Includes feasibility studies which involves financial and legal assessment.
			\item<4-> May take years to complete
		\end{itemize}
	\end{frame}

\subsection*{Design}
	\begin{frame}{Design}
		\begin{itemize}
			\item<2-> Involves two stages: preliminary and final
			\item<3-> Includes the concept, scope, structure, materials, method of construction, and cost and timeline estimate 
		\end{itemize}
	\end{frame}

\subsection*{Construction}
	\begin{frame}{Construction}
		\begin{itemize}
			\item<2-> The actual implementation of the project.
			\item<3-> Includes the monitoring of financial and time aspects of the project, as well as safety and environmental regulations
		\end{itemize}	
\end{frame}

\subsection*{Maintenance/Operation}
	\begin{frame}{Maintenance/Operation}
		\begin{itemize}
			\item<1-> Begins when the construction phase ends, and the owner accepts the completed project.
		\end{itemize}
	\end{frame}

%------------------------------------------------



%------------------------------------------------

%\section*{References}
%\SectionPage

%------------------------------------------------

%\begin{frame} % Use [allowframebreaks] to allow automatic splitting across slides if the content is too long
%	\frametitle{References}
%	
%	\begin{thebibliography}{99} % Beamer does not support BibTeX so references must be inserted manually as below, you may need to use multiple columns and/or reduce the font size further if you have many references
%		\footnotesize % Reduce the font size in the bibliography
%		
%		\bibitem[Smith, 2022]{p1}
%			John Smith (2022)
%			\newblock Publication title
%			\newblock \emph{Journal Name} 12(3), 45 -- 678.
%			
%		\bibitem[Kennedy, 2023]{p2}
%			Annabelle Kennedy (2023)
%			\newblock Publication title
%			\newblock \emph{Journal Name} 12(3), 45 -- 678.
%	\end{thebibliography}
%\end{frame}


%----------------------------------------------------------------------------------------
%	CLOSING SLIDE
%----------------------------------------------------------------------------------------

\begin{frame}[plain] % The optional argument 'plain' hides the headline and footline
	\begin{center}
		{\Huge The End}
		
		\bigskip\bigskip % Vertical whitespace
		
		{\LARGE Questions? Comments?}
	\end{center}
\end{frame}

%----------------------------------------------------------------------------------------

\end{document} 